\documentclass[12pt]{article}
\usepackage[utf8]{inputenc}
\usepackage[T1]{fontenc}
\usepackage[margin=1in]{geometry}
\usepackage{setspace}
\usepackage{xcolor}
\usepackage[colorlinks=true, urlcolor=blue!60!black, linkcolor=blue!60!black]{hyperref}
\setlength{\parindent}{0em}
\setlength{\parskip}{0.6em}
\renewcommand{\baselinestretch}{1.0}
\begin{document}
\begin{center}
{\LARGE \textit{United States v. Texas}, 599 U.S. 670 (2023)\par}
\vspace{0.5em}
{\large Case Brief --- Two-Page Edition\par}
\end{center}

\section{Memory Jogger}
States cannot sue to force the Executive to make more immigration arrests: indirect costs of non-enforcement are not a judicially cognizable injury (\href{https://www.law.cornell.edu/supremecourt/text/22-58}{Cornell LII}).

\section{Detailed Case Facts}
DHS's 2021 Guidelines prioritized arresting noncitizens who threaten national security, public safety, or border security. Texas and Louisiana claimed the Guidelines violated statutes saying DHS ``shall'' take custody of and detain certain noncitizens, and that the States bore resulting supervision, recidivism, education, and healthcare costs.

\section{Procedural History}
After a bench trial, the district court found standing and vacated the Guidelines nationwide; the Fifth Circuit denied a stay. The Supreme Court granted certiorari before judgment.

\section{Judicial Votes}
Reversed, 8--1. Kavanaugh wrote for Roberts, Sotomayor, Kagan, and Jackson; Gorsuch (with Thomas and Barrett) and Barrett (with Gorsuch) concurred in the judgment; Alito dissented.

\section{Holding}
A State that is neither prosecuted nor threatened with prosecution lacks standing to challenge the Executive's arrest and prosecution policies, even if it bears indirect costs. The Court did not decide whether the Guidelines were lawful.

\section{Analysis of Opinions}
\textbf{Majority:} An injury must be judicially cognizable, and no tradition supports orders compelling arrests (\textit{Linda R. S.}); such suits coerce no one, intrude on Article II enforcement discretion, and require second-guessing resource choices. Five possible exceptions remain open. \textbf{Gorsuch:} the injuries were not redressable, since vacatur leaves discretion intact, and universal vacatur is doubtful. \textbf{Barrett:} the majority overread \textit{Linda R. S.} and fused statutory doctrines into Article III. \textbf{Alito:} Texas proved concrete costs, ``shall'' means must, and the ruling revives royal suspending powers.

\section{Examples: Future Applications}
\textbf{Same-Side:} a state challenging DOJ's deprioritizing of marijuana prosecutions lacks standing. \textbf{Opposite-Side:} a program conferring lawful presence and work authorization implicates more than arrest discretion and may support standing. \textbf{Fence-Sitter:} a statute expressly authorizing states to sue to compel immigration custody.

\section{Critique}
Some praised the curb on state-attorney-general suits and Gorsuch's attack on universal vacatur; textualists and Alito's allies say the Court let agencies ignore mandatory statutes, and the \textit{Harvard Law Review} faulted its selective use of history.

\section{Key Quotations}
\begin{quote}
``The States essentially want the Federal Judiciary to order the Executive Branch to alter its arrest policy so as to make more arrests. But this Court has long held `that a citizen lacks standing to contest the policies of the prosecuting authority when he himself is neither prosecuted nor threatened with prosecution.' \textit{Linda R. S. v. Richard D.}, 410 U. S. 614, 619 (1973). Consistent with that fundamental Article III principle, we conclude that the States lack Article III standing to bring this suit.''
\end{quote}
\noindent --- Justice Kavanaugh, for the Court.

\end{document}
